\documentclass[11pt,letterpaper]{article}

% --- PAQUETES Y CONFIGURACIÓN GENERAL ---
\usepackage[utf8]{inputenc}
\usepackage[spanish,es-tabla]{babel}
\usepackage[margin=2.2cm,top=2.5cm,bottom=2.5cm]{geometry}
\usepackage{amsmath,amssymb,amsfonts}
\usepackage{graphicx}
\usepackage{xcolor}
\usepackage{booktabs}
\usepackage{multirow}
\usepackage{tabularx}
\usepackage{longtable}
\usepackage{array}
\usepackage{colortbl}
\usepackage{fancyhdr}
\usepackage{enumitem}
\usepackage{titlesec}
\usepackage{tcolorbox}
\usepackage[hidelinks,pdftex,
  pdftitle={E1.3.1 - Propuesta Preliminar de Equipos y Parámetros de Diseño},
  pdfauthor={Solvanta Engineering Solutions S.A. de C.V.}]{hyperref}

% --- PALETA DE COLORES CORPORATIVA SOLVANTA ---
\definecolor{SolvantaNavy}{HTML}{0B2545}
\definecolor{SolvantaTeal}{HTML}{134074}
\definecolor{SolvantaCyan}{HTML}{8DA9C4}
\definecolor{SolvantaLightGray}{HTML}{F4F7F9}
\definecolor{SolvantaDarkGray}{HTML}{333333}
\definecolor{SolvantaGreen}{HTML}{00A86B}

% --- ESTILO DE ENCABEZADOS Y PIES DE PÁGINA ---
\pagestyle{fancy}
\fancyhf{}
\lhead{\small\textbf{\color{SolvantaNavy}SOLVANTA | Engineering Solutions}}
\rhead{\small\color{SolvantaTeal}RFP-CCM-EM-001 | Entregable E1.3.1 (EQ-01)}
\lfoot{\small\color{gray}Propuesta Preliminar de Equipos y Parámetros de Diseño}
\rfoot{\small\color{gray}Página \thepage}
\renewcommand{\headrulewidth}{0.6pt}
\renewcommand{\footrulewidth}{0.4pt}

% --- FORMATO DE TÍTULOS DE SECCIÓN ---
\titleformat{\section}{\Large\bfseries\color{SolvantaNavy}}{\thesection.}{0.5em}{}[\titlerule]
\titleformat{\subsection}{\large\bfseries\color{SolvantaTeal}}{\thesubsection.}{0.5em}{}
\titleformat{\subsubsection}{\bfseries\color{SolvantaNavy}}{\thesubsubsection.}{0.5em}{}

% --- COMANDOS ESPECIALES ---
\newcommand{\tagbox}[1]{\textbf{\color{SolvantaNavy}#1}}
\newcolumntype{Y}{>{\centering\arraybackslash}X}

\begin{document}

% ==============================================================================
% PORTADA Y ENCABEZADO TÉCNICO
% ==============================================================================
\begin{titlepage}
    \centering
    \vspace*{0.5cm}
    {\huge \textbf{\color{SolvantaNavy}SOLVANTA ENGINEERING SOLUTIONS S.A. DE C.V.}}\\[0.4cm]
    {\large \textbf{\color{SolvantaTeal}INGENIERÍA DE PROCESOS \& SUSTENTABILIDAD INDUSTRIAL}}\\[1.5cm]

    \begin{tcolorbox}[colback=SolvantaLightGray,colframe=SolvantaNavy,arc=3mm,boxrule=1.2pt]
        \centering \vspace{0.2cm}
        {\Large \textbf{\color{SolvantaNavy}LICITACIÓN PÚBLICA INTERNACIONAL RFP-CCM-EM-001}}\\[0.2cm]
        {\large \textbf{\color{SolvantaTeal}PLANTA DE PRODUCCIÓN DE E-METANOL VERDE (100,000 t/año)}}\\[0.3cm]
        \rule{\linewidth}{0.5pt}
        \vspace{0.2cm}
        {\LARGE \textbf{\color{SolvantaNavy}ENTREGABLE E1.3.1 (ACTIVIDAD EQ-01)}}\\[0.2cm]
        {\Large \textbf{\color{SolvantaTeal}Propuesta Preliminar de Equipos y Parámetros de Diseño previa a DWSIM v0}}\\[0.2cm]
    \end{tcolorbox}

    \vspace{1.5cm}

    \begin{table}[h!]
        \centering
        \small
        \begin{tabularx}{0.95\textwidth}{|l|X|c|}
            \hline
            \rowcolor{SolvantaNavy}
            \textbf{\color{white}Rol RACI} & \textbf{\color{white}Integrante / Identidad} & \textbf{\color{white}Usuario GitHub} \\
            \hline
            \textbf{Responsable (R) \& Autor} & Karol Paola Camacho López (Ingeniera Principal de Procesos) & \texttt{@Karoc} \\
            \hline
            \textbf{Aprobador / Consultado (C)} & Gildardo Ocaña Barrales (Ingeniero de Diagramación y 3D) & \texttt{@GildardoOca25} \\
            \hline
            \textbf{Informado (I)} & Paulina Mota Macip (Ingeniera de Simulación) & \texttt{@paulina145} \\
            \hline
            \textbf{Informada (I)} & Grecia Figueroa Rincón (Ingeniera de Normativa \& HSE) & \texttt{@Cho0bb} \\
            \hline
        \end{tabularx}
    \end{table}

    \vfill

    \begin{tcolorbox}[colback=white,colframe=SolvantaTeal,arc=2mm,boxrule=0.8pt]
        \small \centering
        \textbf{Documento de Ingeniería Oficial — Solvanta Engineering Solutions S.A. de C.V.}\\[0.1cm]
        Fecha de Emisión: 27 de Agosto de 2026 \quad | \quad Versión: 1.0 \quad | \quad Estado: Aprobado para Simulación DWSIM v0
    \end{tcolorbox}
\end{titlepage}

\newpage

% ==============================================================================
% 1. PROPÓSITO DE LA ACTIVIDAD Y PRINCIPIO DE CONCORDANCIA
% ==============================================================================
\section{Propósito de la Actividad y Principio de Concordancia}

El presente documento constituye el **Entregable E1.3.1 (Actividad Complementaria EQ-01)**, elaborado por la firma de ingeniería **Solvanta Engineering Solutions S.A. de C.V.** como puente analítico indispensable entre el *Diagrama de Bloques del Proceso (BFD v0 - Entregable E1.3)* y la construcción del modelo riguroso de simulación termodinámica en *DWSIM v0 (Entregable E1.2)*.

Hasta la etapa E1.3, la firma ha definido rigurosamente **qué funciones** debe realizar la planta de e-metanol verde. El objetivo prioritario de la actividad EQ-01 consiste en definir **qué equipos u operaciones unitarias** realizarán cada una de esas funciones, traduciendo la arquitectura conceptual del BFD v0 en una propuesta tecnológica preliminar sustentada en parámetros de diseño.

\begin{tcolorbox}[colback=SolvantaLightGray,colframe=SolvantaNavy,title=\textbf{Principio de Concordancia Fundamental (Regla BFD $\rightarrow$ EQ-01)}]
\small
Todo equipo propuesto en el presente documento responde estrictamente a la pregunta:
\begin{center}
    \textit{¿Qué bloque o necesidad técnica del BFD v0 justifica que este equipo exista?}
\end{center}
Ningún equipo ha sido incorporado por inercia técnica o costumbre; cada unidad u operación unitaria cuenta con un TAG normalizado, corrientes de entrada y salida asociadas, y una justificación funcional trazable.
\end{tcolorbox}

---

% ==============================================================================
% 2. MATRIZ BLOQUE - EQUIPO (SECCIÓN 3 DE EQ-01)
% ==============================================================================
\section{Matriz Bloque -- Equipo (Traducción del BFD v0 a Operaciones Unitarias)}

A continuación se presenta la matriz completa que vincula cada bloque funcional del BFD v0 con las operaciones unitarias y equipos preliminares requeridos para la planta de e-metanol (100,000 t/año).

\begin{longtable}{|p{1.4cm}|p{3.2cm}|p{3.0cm}|p{1.6cm}|p{4.8cm}|}
    \hline
    \rowcolor{SolvantaNavy}
    \textbf{\color{white}Bloque BFD} & \textbf{\color{white}Función del Bloque} & \textbf{\color{white}Equipo Propuesto} & \textbf{\color{white}TAG Prelim.} & \textbf{\color{white}Justificación Técnica Preliminar} \\
    \hline
    \endfirsthead
    \hline
    \rowcolor{SolvantaNavy}
    \textbf{\color{white}Bloque BFD} & \textbf{\color{white}Función del Bloque} & \textbf{\color{white}Equipo Propuesto} & \textbf{\color{white}TAG Prelim.} & \textbf{\color{white}Justificación Técnica Preliminar} \\
    \hline
    \endhead

    \multirow{3}{*}{\textbf{B-100}} & \multirow{3}{=}{Recepción y acondicionamiento de $\text{CO}_2$ de captura externa} & Compresor de $\text{CO}_2$ & \tagbox{C-101 A/B} & Eleva la presión del $\text{CO}_2$ desde 1.5 bar de succión hasta 20 bar. \\ \cline{3-5}
     & & Enfriador Interapa & \tagbox{E-101} & Condensa la humedad sobrante post-compresión previa al SOEC. \\ \cline{3-5}
     & & Separador de Humedad & \tagbox{V-101} & Remueve condensados acuosos para proteger el compresor y celda. \\ \hline

    \multirow{3}{*}{\textbf{B-200}} & \multirow{3}{=}{Co-electrólisis SOEC / Generación de Syngas a $750^\circ\text{C}$} & Sobrecalentador de Reactivos & \tagbox{E-201} & Eleva la temperatura de vapor y $\text{CO}_2$ de $200^\circ\text{C}$ a $750^\circ\text{C}$ con integración. \\ \cline{3-5}
     & & Stacks de Electrólisis SOEC & \tagbox{R-201} & Co-electrólisis catódica para disociar $\text{H}_2\text{O}$ y $\text{CO}_2$ generando Syngas. \\ \cline{3-5}
     & & Separador de $\text{O}_2$ Anódico & \tagbox{V-201} & Acondiciona la corriente de $\text{O}_2$ puro producido como coproducto. \\ \hline

    \multirow{2}{*}{\textbf{B-300}} & \multirow{2}{=}{Compresión de Gas de Síntesis (Syngas)} & Compresor Principal de Syngas & \tagbox{C-301 A/B} & Compresión multietapa del syngas desde 20 bar hasta 75 bar de reacción. \\ \cline{3-5}
     & & Enfriador Interetapas & \tagbox{E-301} & Reduce la temperatura de descarga y optimiza potencia de compresión. \\ \hline

    \textbf{B-400} & Síntesis Catalítica de Metanol & Reactor Multitubular de Metanol & \tagbox{R-401} & Reactor exotérmico a $240^\circ\text{C}$ y 75 bar sobre catalizador $\text{CuO/ZnO/Al}_2\text{O}_3$. Genera vapor en carcasa. \\ \hline

    \multirow{2}{*}{\textbf{B-500}} & \multirow{2}{=}{Separación Primaria y Condensación Flash} & Condensador de Efluente & \tagbox{E-501} & Enfría el efluente del reactor de $240^\circ\text{C}$ a $38^\circ\text{C}$ para condensar metanol y agua. \\ \cline{3-5}
     & & Separador Flash de Alta Presión & \tagbox{V-501} & Separa líquido condensado (metanol crudo) del syngas no reaccionado. \\ \hline

    \multirow{2}{*}{\textbf{B-600}} & \multirow{2}{=}{Lazo de Reciclo y Purga Inteligente} & Divisor de Purga de Inertes & \tagbox{SD-601} & Expulsa purga del 2.5\% para evitar acumulación de $\text{N}_2$ en el lazo. \\ \cline{3-5}
     & & Compresor de Reciclo & \tagbox{C-601} & Recomprime el gas no reaccionado venconendos caídas de presión ($\Delta P \approx 3\text{ bar}$). \\ \hline

    \multirow{4}{*}{\textbf{B-700}} & \multirow{4}{=}{Purificación por Destilación Fraccionada} & Columna Fraccionadora & \tagbox{T-701} & Separación de metanol refinado ($\ge 99.5\%$ wt) de agua residual (24 platos). \\ \cline{3-5}
     & & Condensador de Reflujo & \tagbox{E-702} & Condensa totalmente los vapores de metanol en la cabeza de columna. \\ \cline{3-5}
     & & Rehervidor de Fondo & \tagbox{E-701} & Suministra calor de vaporización usando vapor de proceso recuperado. \\ \cline{3-5}
     & & Acumulador de Reflujo & \tagbox{V-701} & Recipiente de separación de metanol líquido refinado y retorno. \\ \hline

    \multirow{2}{*}{\textbf{B-800}} & \multirow{2}{=}{Almacenamiento y Despacho Comercial} & Bomba de Transferencia & \tagbox{P-801 A/B} & Bombea el metanol refinado a tanques de almacenamiento atmosférico. \\ \cline{3-5}
     & & Tanque de Almacenamiento & \tagbox{TK-801} & Almacenamiento seguro de producto refinado bajo colchón de nitrógeno. \\ \hline
\end{longtable}

---

% ==============================================================================
% 3. MATRIZ DE PARÁMETROS DE DISEÑO POR DETERMINAR (SECCIÓN 4 DE EQ-01)
% ==============================================================================
\section{Matriz de Parámetros de Diseño por Determinar}

Para cada equipo propuesto, se definen a continuación los parámetros críticos de diseño, su unidad de medida, su clasificación de estado actual (siguiendo estrictamente las categorías de EQ-01) y la fuente/evidencia futura prevista.

\begin{small}
\begin{longtable}{|p{1.2cm}|p{2.2cm}|p{3.2cm}|p{1.2cm}|p{2.4cm}|p{2.6cm}|p{2.4cm}|}
    \hline
    \rowcolor{SolvantaNavy}
    \textbf{\color{white}TAG} & \textbf{\color{white}Equipo} & \textbf{\color{white}Parámetro por Determinar} & \textbf{\color{white}Unidad} & \textbf{\color{white}Estado} & \textbf{\color{white}Método Previsto} & \textbf{\color{white}Fuente Futura} \\
    \hline
    \endfirsthead
    \hline
    \rowcolor{SolvantaNavy}
    \textbf{\color{white}TAG} & \textbf{\color{white}Equipo} & \textbf{\color{white}Parámetro por Determinar} & \textbf{\color{white}Unidad} & \textbf{\color{white}Estado} & \textbf{\color{white}Método Previsto} & \textbf{\color{white}Fuente Futura} \\
    \hline
    \endhead

    \textbf{C-101} & Compresor $\text{CO}_2$ & Potencia de compresión & kW & Por simular & Modelo politrópico & DWSIM v0 \\ \hline
    \textbf{C-101} & Compresor $\text{CO}_2$ & Presión de descarga & bar & Supuesto temporal & Balance de succión & Base de diseño \\ \hline
    \textbf{E-201} & Sobrecalentador & Carga térmica ($\dot{Q}$) & kW & Por simular & Balance de entalpía & DWSIM / BioSTEAM \\ \hline
    \textbf{R-201} & Stack SOEC & Consumo eléctrico específico & $\frac{\text{kWh}}{\text{kg }\text{H}_2}$ & Dato disponible & Medición de referencia & RFP A2.1 / NB-01 \\ \hline
    \textbf{R-201} & Stack SOEC & Eficiencia Faraday & \% & Supuesto temporal & Modelo electroquímico & Literatura SOEC \\ \hline
    \textbf{C-301} & Compresor Syngas & Eficiencia isentrópica & \% & Supuesto temporal & Correlación de diseño & Catálogo fabricante \\ \hline
    \textbf{C-301} & Compresor Syngas & Número de etapas & -- & Por calcular & Relación $P_2/P_1 \le 3.5$ & Cálculo en Python \\ \hline
    \textbf{R-401} & Reactor Metanol & Conversión por paso $\text{CO}_2$ & \% & Por simular & Equilibrio / Cinética & DWSIM Kinetic Model \\ \hline
    \textbf{R-401} & Reactor Metanol & Masa de catalizador & kg & Por calcular & Tiempo espacial / WHSV & Literatura Lurgi \\ \hline
    \textbf{R-401} & Reactor Metanol & Vapor generado en carcasa & t/h & Por simular & Balance de calor & DWSIM v0 \\ \hline
    \textbf{E-501} & Enfriador Flash & Temperatura de descarga & $^\circ\text{C}$ & Dato disponible & Especificación de proceso & NB-01 \\ \hline
    \textbf{V-501} & Separador Flash & Fracción vaporizada & -- & Por simular & Flash PT Peng-Robinson & DWSIM v0 \\ \hline
    \textbf{SD-601} & Divisor Purga & Porcentaje de purga & \% & Por optimizar & Minimización de inertes & Simulación Python \\ \hline
    \textbf{C-601} & Compresor Reciclo & Caudal de reciclo volumétrico & $\text{m}^3/\text{h}$ & Por simular & Balance del lazo cerrado & DWSIM v0 \\ \hline
    \textbf{T-701} & Columna Destilac. & Número de platos teóricos & -- & Por calcular & Método Fenske-Underwood & DWSIM Shortcut \\ \hline
    \textbf{T-701} & Columna Destilac. & Relación de reflujo ($R/R_{\min}$) & -- & Por optimizar & Análisis costo rehervidor & DWSIM Rigorous \\ \hline
    \textbf{T-701} & Columna Destilac. & Pureza de metanol comercial & \% wt & Dato disponible & Especificación cliente & RFP $\ge 99.5\%$ \\ \hline
\end{longtable}
\end{small}

---

% ==============================================================================
% 4. FICHAS PRELIMINARES DE EQUIPOS CRÍTICOS (AL MENOS 5 FICHAS)
% ==============================================================================
\section{Fichas Preliminares por Equipo Crítico}

Se seleccionan cinco equipos estratégicos clave para el proceso global de e-metanol verde.

\subsection{Ficha Crítica 1: Reactor Catalítico Multitubular de Metanol (R-401)}
\begin{tcolorbox}[colback=white,colframe=SolvantaNavy,title=\textbf{Ficha Técnica Preliminar — R-401}]
\small
\begin{description}[font=\bfseries\color{SolvantaTeal}]
    \item[TAG Preliminar:] R-401
    \item[Nombre del Equipo:] Reactor Catalítico Multitubular de Síntesis de Metanol (Tipo Lurgi)
    \item[Bloque BFD Asociado:] B-400 (Síntesis de Metanol)
    \item[Función Principal:] Llevar a cabo la conversión exotérmica catalítica de syngas ($\text{CO}_2 + 3\text{H}_2 \rightleftharpoons \text{CH}_3\text{OH} + \text{H}_2\text{O}$) a $240^\circ\text{C}$ y 75.0 bar.
    \item[Corrientes de Entrada:] Syngas comprimido mezclado con reciclo (C-06 + C-09b).
    \item[Corrientes de Salida:] Efluente de síntesis de metanol crudo a $240^\circ\text{C}$ (C-07).
    \item[Parámetros por Determinar:] Conversión por paso de $\text{CO}_2$, perfil de temperatura radial y axial, caída de presión en el lecho catalítico, masa de catalizador $\text{CuO/ZnO/Al}_2\text{O}_3$ requerida.
    \item[Información esperada de DWSIM v0:] Conversión de equilibrio cinético, composición molar exacta del efluente y flujo másico de vapor de agua a $210^\circ\text{C}$ generado en el envolvente.
    \item[Principal Incertidumbre Técnica:] Desactivación del catalizador por envenenamiento con trazas de azufre o sinterización térmica por puntos calientes (*hotspots*).
\end{description}
\end{tcolorbox}

\subsection{Ficha Crítica 2: Stacks de Electrólisis SOEC de Alta Temperatura (R-201)}
\begin{tcolorbox}[colback=white,colframe=SolvantaNavy,title=\textbf{Ficha Técnica Preliminar — R-201}]
\small
\begin{description}[font=\bfseries\color{SolvantaTeal}]
    \item[TAG Preliminar:] R-201
    \item[Nombre del Equipo:] Unidad de Stacks de Co-Electrólisis de Óxido Sólido (SOEC)
    \item[Bloque BFD Asociado:] B-200 (Generación de Syngas por SOEC)
    \item[Función Principal:] Co-electrolizar vapor de agua sobrecalentado y $\text{CO}_2$ a $750^\circ\text{C}$ para producir gas de síntesis directamente con la relación $\text{H}_2/\text{CO}$ deseada.
    \item[Corrientes de Entrada:] Mezcla de vapor de agua y $\text{CO}_2$ sobrecalentado a $750^\circ\text{C}$ (C-01b + C-02), Energía eléctrica y térmica.
    \item[Corrientes de Salida:] Gas de síntesis catódico (C-04) y Oxígeno puro anódico (C-05).
    \item[Parámetros por Determinar:] Consumo eléctrico específico neto ($\approx 41\text{ kWh/kg }\text{H}_2$), degradación celular (\%/1000h), relación área/potencia.
    \item[Información esperada de DWSIM v0:] Balance de materia y energía simplificado considerando estequiometría de disociación e integración de calor residual.
    \item[Principal Incertidumbre Técnica:] Integridad mecánica de los sellos cerámicos bajo ciclos térmicos severos a $750^\circ\text{C}$.
\end{description}
\end{tcolorbox}

\subsection{Ficha Crítica 3: Compresor Multietapa de Gas de Síntesis (C-301 A/B)}
\begin{tcolorbox}[colback=white,colframe=SolvantaNavy,title=\textbf{Ficha Técnica Preliminar — C-301 A/B}]
\small
\begin{description}[font=\bfseries\color{SolvantaTeal}]
    \item[TAG Preliminar:] C-301 A/B (Arreglo 2 x 50\%)
    \item[Nombre del Equipo:] Compresor Centrífugo Multietapa de Syngas
    \item[Bloque BFD Asociado:] B-300 (Compresión de Syngas)
    \item[Función Principal:] Elevar la presión del gas de síntesis húmedo desde la presión de salida del SOEC ($\approx 20\text{ bar}$) hasta la presión del lazo de síntesis ($75.0\text{ bar}$).
    \item[Corrientes de Entrada:] Syngas frío acondicionado a 20 bar (C-04).
    \item[Corrientes de Salida:] Syngas comprimido a 75 bar (C-06).
    \item[Parámetros por Determinar:] Potencia de accionamiento eléctrica (kW), número exacto de etapas con enfriamiento intermedio, temperatura máxima de descarga.
    \item[Información esperada de DWSIM v0:] Potencia isentrópica y consumo de agua de enfriamiento en los intercambiadores interetapa.
    \item[Principal Incertidumbre Técnica:] Fragilización por hidrógeno en aleaciones de los álabes del compresor a 75 bar.
\end{description}
\end{tcolorbox}

\subsection{Ficha Crítica 4: Separador Flash de Alta Presión (V-501)}
\begin{tcolorbox}[colback=white,colframe=SolvantaNavy,title=\textbf{Ficha Técnica Preliminar — V-501}]
\small
\begin{description}[font=\bfseries\color{SolvantaTeal}]
    \item[TAG Preliminar:] V-501
    \item[Nombre del Equipo:] Separador Trifásico / Flash de Alta Presión
    \item[Bloque BFD Asociado:] B-500 (Separación Primaria)
    \item[Función Principal:] Realizar la separación líquido-vapor del efluente enfriado a $38^\circ\text{C}$ y 73 bar, obteniendo metanol crudo líquido en el fondo y reciclo gaseoso por el domo.
    \item[Corrientes de Entrada:] Efluente de síntesis bifásico enfriado (C-07b).
    \item[Corrientes de Salida:] Vapor no reaccionado hacia lazo de reciclo (C-09) y Metanol crudo líquido (C-08).
    \item[Parámetros por Determinar:] Dimensiones del recipiente (diámetro y altura para evitar arrastre de gotas), tiempo de residencia del líquido ($\approx 5\text{ min}$).
    \item[Información esperada de DWSIM v0:] Equilibrio líquido-vapor (LVE) calculado mediante la ecuación de estado Peng-Robinson y solubilidad de gases desgasificados.
    \item[Principal Incertidumbre Técnica:] Pérdida de metanol en fase vapor por arrastre de neblina a elevada presión.
\end{description}
\end{tcolorbox}

\subsection{Ficha Crítica 5: Columna de Destilación Fraccionada de Metanol (T-701)}
\begin{tcolorbox}[colback=white,colframe=SolvantaNavy,title=\textbf{Ficha Técnica Preliminar — T-701}]
\small
\begin{description}[font=\bfseries\color{SolvantaTeal}]
    \item[TAG Preliminar:] T-701
    \item[Nombre del Equipo:] Columna de Destilación Fraccionada de Metanol (24 Platos Teóricos)
    \item[Bloque BFD Asociado:] B-700 (Purificación de Metanol)
    \item[Función Principal:] Purificar el metanol crudo líquido desde $\approx 70\%$ wt hasta cumplir la especificación de pureza comercial $\ge 99.5\%$ wt (Grado AA).
    \item[Corrientes de Entrada:] Alimentación de metanol crudo desgasificado.
    \item[Corrientes de Salida:] Metanol producto por domo ($\ge 99.5\%$ wt) y Agua purificada de fondo enviada a reciclo ZLD.
    \item[Parámetros por Determinar:] Número real de platos, plato óptimo de alimentación, relación de reflujo de operación, carga térmica del rehervidor y condensador.
    \item[Información esperada de DWSIM v0:] Perfil de temperaturas plato a plato y perfil de composiciones mediante simulación rigurosa RadFrac / Rigorous Column.
    \item[Principal Incertidumbre Técnica:] Formación de azeótropos secundarios o acumulación de alcoholes superiores (fusel oil) en platos intermedios.
\end{description}
\end{tcolorbox}

---

% ==============================================================================
% 5. TABLA DE CONCORDANCIA CON DWSIM V0 (SECCIÓN 8 DE EQ-01)
% ==============================================================================
\section{Tabla de Concordancia con DWSIM v0}

Previo a la construcción del *flowsheet* de simulación en DWSIM v0, la siguiente tabla establece la equivalencia precisa de bloques entre la propuesta física de equipos y los módulos de simulación disponibles en la plataforma.

\begin{table}[h!]
    \centering
    \small
    \begin{tabularx}{\textwidth}{|l|X|l|X|}
        \hline
        \rowcolor{SolvantaNavy}
        \textbf{\color{white}TAG} & \textbf{\color{white}Equipo Propuesto} & \textbf{\color{white}Módulo DWSIM} & \textbf{\color{white}Dato a Validar u Obtener} \\
        \hline
        \textbf{C-101} & Compresor de $\text{CO}_2$ & \texttt{Compressor} & Potencia eléctrica y temperatura de descarga \\ \hline
        \textbf{E-101} & Enfriador Interapa & \texttt{Cooler} & Carga térmica y masa de agua condensada \\ \hline
        \textbf{V-101} & Knock-Out Drum & \texttt{Gas-Liquid Separator} & Fracción de líquido separado \\ \hline
        \textbf{E-201} & Sobrecalentador & \texttt{Heater} & Demanda entálpica a $750^\circ\text{C}$ \\ \hline
        \textbf{R-201} & Celda SOEC & \texttt{Conversion Reactor} & Estequiometría de disociación y consumo $\text{kWh/kg}$ \\ \hline
        \textbf{C-301} & Compresor Syngas & \texttt{Compressor} & Potencia global de compresión a 75 bar \\ \hline
        \textbf{R-401} & Reactor de Metanol & \texttt{Equilibrium / PFR} & Conversión cinético-termodinámica y calor generado \\ \hline
        \textbf{E-501} & Enfriador Efluente & \texttt{Cooler} & Temperatura de rocío/burbuja \\ \hline
        \textbf{V-501} & Separador Flash & \texttt{Flash Vessel (PT)} & Composición del líquido crudo y gas reciclo \\ \hline
        \textbf{SD-601} & Divisor Purga & \texttt{Splitter} & Balance de reciclo y acumulación de $\text{N}_2$ \\ \hline
        \textbf{C-601} & Compresor Reciclo & \texttt{Compressor} & Potencia para vencer $\Delta P$ del lazo \\ \hline
        \textbf{T-701} & Columna Destilada & \texttt{Shortcut / Rigorous Column} & Número de platos y relación de reflujo $R/R_{\min}$ \\ \hline
    \end{tabularx}
\end{table}

---

% ==============================================================================
% 6. PREGUNTAS DE ANÁLISIS RIGUROSO (SECCIÓN 9 DE EQ-01)
% ==============================================================================
\section{Preguntas de Análisis Riguroso}

\begin{enumerate}[label=\textbf{\arabic*.}]
    \item \textbf{¿Qué bloque del BFD requirió el mayor número de equipos preliminares y por qué?}\\[0.1cm]
    \textit{Respuesta:} El **Bloque B-700 (Purificación de Metanol por Destilación Fraccionada)** y el **Bloque B-100/B-300 (Acondicionamiento y Compresión)** requirieron el mayor número de equipos individuales. Esto se debe a que la purificación no es una operación unitaria simple, sino un subsistema compuesto por la columna de platos (\tagbox{T-701}), el rehervidor de fondo (\tagbox{E-701}), el condensador de cabezas (\tagbox{E-702}), el acumulador de reflujo (\tagbox{V-701}) y bombas de recirculación (\tagbox{P-801}). La necesidad de alcanzar la estricta pureza comercial de $\ge 99.5\%$ wt exigida por el RFP requiere control de temperatura y transferencia de masa en múltiple etapas.

    \item \textbf{¿Qué equipo consideran más crítico para alcanzar la producción de metanol especificada?}\\[0.1cm]
    \textit{Respuesta:} El **Reactor Catalítico Multitubular de Metanol (\tagbox{R-401})**. Este equipo es el corazón del proceso químico; de su desempeño depende la conversión por paso de $\text{CO}_2$ e $\text{H}_2$, la selectividad hacia metanol versus subproductos (como dimetil éter o alcoholes superiores) y la masa total de metanol crudo generada por hora para garantizar las 100,000 t/año.

    \item \textbf{¿Qué equipo consideran más crítico desde el punto de vista energético?}\\[0.1cm]
    \textit{Respuesta:} La **Unidad Stacks de Electrólisis SOEC (\tagbox{R-201})** (o en el Caso Base, la planta PEM). Concentra más del 85\% del consumo de energía eléctrica del complejo industrial ($\approx 41.0\text{ kWh/kg }\text{H}_2$). Asimismo, desde el punto de vista térmico, la integración del reactor \tagbox{R-401} (que genera vapor a $210^\circ\text{C}$) con el precalentador \tagbox{E-201} es determinante para evitar consumo de combustible externo.

    \item \textbf{¿Qué parámetros ya pueden fijarse a partir de las Bases de Diseño y NB-01?}\\[0.1cm]
    \textit{Respuesta:}
    \begin{itemize}[leftmargin=*]
        \item Capacidad nominal de producción de e-metanol: $100,000\text{ t/año}$ ($12.5\text{ t/h}$ a 8,000 h/año).
        \item Pureza del producto final comercial: $\ge 99.5\%$ en masa.
        \item Presión del lazo de síntesis de metanol: $75.0\text{ bar}$.
        \item Temperatura de operación del reactor catalítico: $240^\circ\text{C}$.
        \item Relación estequiométrica ideal de alimentación al reactor: $\text{H}_2/\text{CO}_2 = 3.0$.
        \item Temperatura del agua de enfriamiento disponible en batería: $25^\circ\text{C}$.
    \end{itemize}

    \item \textbf{¿Qué parámetros no pueden determinarse todavía sin utilizar DWSIM?}\\[0.1cm]
    \textit{Respuesta:}
    \begin{itemize}[leftmargin=*]
        \item La composición molar exacta de equilibrio y la velocidad de reciclo en el lazo cerrado B-500/B-600.
        \item La fracción óptima de purga de gases inertes ($\text{N}_2$) en el divisor \tagbox{SD-601}.
        \item Las cargas térmicas exactas ($\dot{Q}$) del condensador efluente \tagbox{E-501} y rehervidor \tagbox{E-701}.
        \item El número óptimo de platos reales y la relación de reflujo en la columna de destilación \tagbox{T-701}.
    \end{itemize}

    \item \textbf{¿Qué decisiones podrían provocar posteriormente una modificación del BFD v0?}\\[0.1cm]
    \textit{Respuesta:}
    \begin{enumerate}[label=\alph*)]
        \item **Resultados del Modelo Tecno-Económico PEM vs. SOEC:** Si la simulación rigurosa demuestra que el balance térmico SOEC a $750^\circ\text{C}$ requiere una fuente auxiliar eléctrica excesiva que opaca el ahorro de OPEX, se ratificaría el BFD del Caso Base PEM.
        \item **Estrategia de Purga e Inertes:** Si la acumulación de $\text{N}_2$ en el lazo de reciclo reduce drásticamente la presión parcial de reactivos en el reactor \tagbox{R-401}, podría ser necesario agregar un bloque adicional de separación de gases con membrana o PSA en el reciclo.
    \end{enumerate}
\end{enumerate}

---

% ==============================================================================
% 7. CONCLUSIÓN Y LISTA DE COMPROBACIÓN
% ==============================================================================
\section{Conclusión y Estado de Preparación para DWSIM v0}

Con la conclusión del **Entregable E1.3.1 (EQ-01)**, la firma **Solvanta Engineering Solutions S.A. de C.V.** ha completado la traducción estructural del BFD v0 a un esquema técnico de equipos preliminares con TAGs normalizados y parámetros clasificados. Este trabajo proporciona la base física precisa para la construcción del modelo en **DWSIM v0**.

\begin{table}[h!]
    \centering
    \small
    \begin{tabularx}{0.9\textwidth}{|X|c|}
        \hline
        \rowcolor{SolvantaNavy}
        \textbf{\color{white}Criterio de Verificación de Calidad (EQ-01)} & \textbf{\color{white}Estado} \\
        \hline
        Cada equipo se relaciona con un bloque del BFD v0 & \textbf{Cumplido (100\%)} \\
        Cada equipo tiene una función claramente definida y justificada & \textbf{Cumplido (100\%)} \\
        Se asignó un TAG preliminar a cada equipo & \textbf{Cumplido (100\%)} \\
        Se identificaron los parámetros por determinar y sus unidades & \textbf{Cumplido (100\%)} \\
        Cada parámetro tiene un estado definido (siguiendo guía de 6 categorías) & \textbf{Cumplido (100\%)} \\
        Se completaron al menos cinco fichas de equipos críticos & \textbf{Cumplido (5/5)} \\
        Existe concordancia técnica directa con los módulos de DWSIM v0 & \textbf{Cumplido (100\%)} \\
        \hline
    \end{tabularx}
\end{table}

\vspace{0.5cm}

\begin{tcolorbox}[colback=SolvantaNavy,colframe=SolvantaNavy]
    \color{white}
    \centering
    \textbf{Firma y Autorización Técnica:}\\
    \vspace{0.2cm}
    \textit{Karol Paola Camacho López (\texttt{@Karoc}) — Ingeniera Principal de Procesos}\\
    \textit{Gildardo Ocaña Barrales (\texttt{@GildardoOca25}) — Ingeniero de Revisión y Consentimiento}
\end{tcolorbox}

\end{document}
